\lecture{8}{16. Marts 2026}{Linear isotropic materials and the uniaxial stress}

\section{Material Behavior}

\subsection{Uniaxial Behavior}
The most elementary and well-known description of material behavior is \textit{Hooke's law}:
\begin{equation*}
  \sigma_{11} = E \epsilon_{11}
\end{equation*}
where $E$ is the modulus of elasticity of Young's modulus.
This also gives rise to stress strain curves as the one shown on \Cref{fig:f4_1}.

\begin{figure} [ht]
  \centering
  \includegraphics[width=0.5\linewidth]{./figures/f4_1.png}
  \caption{Stress-strain curve for two cycles of loading and unloading for mild steel.}
  \label{fig:f4_1}
\end{figure}

Branch (1) is the initial elastic response, from which $E$ may be calculated as the slope.
In the vicinity of (2), there is a decrease in the slope, commencing at the \textit{proportional limit} and progressing until the \textit{yield point} is reached at (3).
Yielding progresses along (4) until unloading commences at (5) and continues along (6).
The unloading ceases and reloading begins at (7) to initiate another cycle that extends through (13) when the test is terminated.
Region (14) is known as a \textit{hysteresis loop} and is a measure of the energy dissipated through one excursion beyond the yield point.

\subsection{Generalized Hooke's Law}
We will now introduce the \textit{strain energy density} $W$ that is a homogeneous quadratic function of the strain components $W(\epsilon_{ij})$.
For a body that is slightly strained by gradual application of the loading while the temperature remains constant, a stress component will be produced derivable as
\begin{equation*}
  \sigma_{ij} = \frac{\partial W}{\partial \epsilon_{ij}} 
.\end{equation*}
Here, it is convenient to write $W$ in terms of the engineering strains $ \overline{\epsilon}_{ij}$.
Thus
\begin{equation*}
  \sigma_{ij} = \frac{\partial W \left( \overline{\epsilon}_{ij} \right)}{\partial \overline{\epsilon}_{ij}} 
\end{equation*}
where $\overline{\epsilon}_{ij} = \left( 2 - \delta_{ij} \right) \epsilon_{ij}$.

The function should have coefficients such that $W \geq 0$ in order to ensure the stability of the body, with $W(0) = 0$ corresponding to a \textit{natural of zero energy state}.

Now, the generalized Hooke's law is written in the form of a fourth-order tensor
\begin{equation*}
  \sigma_{ij} = E_{ijkl} \epsilon_{kl}
\end{equation*}
in which the 81 coefficients $E_{ijkl}$ are called the elastic constants.
Since $W$ is continuous the order of differentiation does not matter and $E_{ijkl}$ is symmetric.
Thus the number of independent equations is reduced from nine to six, the elastic constants from 81 to 36, and then further to $\left( \left( 36 - 6 \right) / 2 \right) + 6 = 21$, when only half of the off-diagonal constants are counted.
We represent this relationship in matrix form in terms of the engineering strains as
\begin{equation*}
  \Vec{\sigma} = \Vec{E} \overline{\Vec{\epsilon}}
\end{equation*}
where
\begin{equation*}
  \left\{ \begin{array}{cccccc}
    \sigma_{11}\\
  \sigma_{22}\\
  \sigma_{33}\\
  \sigma_{12}\\
  \sigma_{23}\\
  \sigma_{31}\\
  \end{array} \right\} = \begin{bmatrix}
    E_{1111} & E_{1122} & E_{1133} & E_{1112} & E_{1123} & E_{1131}\\
   & E_{2222} & E_{2233} & E_{2212} & E_{2223} & E_{2231}\\
   &  & E_{3333} & E_{3312} & E_{3323} & E_{3331}\\
   &  &  & E_{1212} & E_{1223} & E_{1231}\\
   &  &  &  & E_{2323} & E_{2331}\\
   &  &  &  &  & E_{3131}\\
  \end{bmatrix} \left\{ \begin{array}{cccccc}
    \overline{\epsilon}_{11}\\
  \overline{\epsilon}_{22}\\
  \overline{\epsilon}_{33}\\
  \overline{\epsilon}_{12}\\
  \overline{\epsilon}_{23}\\
  \overline{\epsilon}_{31}\\
  \end{array} \right\}
.\end{equation*}
Which can also be written as
\begin{equation*}
  \left\{ \begin{array}{cccccc}
    \sigma_{11}\\
  \sigma_{22}\\
  \sigma_{33}\\
  \sigma_{12}\\
  \sigma_{23}\\
  \sigma_{31}\\
  \end{array} \right\} = \begin{bmatrix}
    C_{11} & C_{12} & C_{13} & C_{14} & C_{15} & C_{16}\\
   & C_{22} & C_{23} & C_{24} & C_{25} & C_{26}\\
   &  & C_{33} & C_{34} & C_{35} & C_{36}\\
   &  &  & C_{44} & C_{45} & C_{46}\\
   &  &  &  & C_{55} & C_{56}\\
   &  &  &  &  &C_{66} \\
  \end{bmatrix} \left\{ \begin{array}{cccccc}
    \overline{\epsilon}_{11}\\
  \overline{\epsilon}_{22}\\
  \overline{\epsilon}_{33}\\
  \overline{\epsilon}_{12}\\
  \overline{\epsilon}_{23}\\
  \overline{\epsilon}_{31}\\
  \end{array} \right\}
\end{equation*}
or
\begin{equation*}
  \Vec{\sigma} =\Vec{C} \overline{\epsilon}
.\end{equation*}
The corresponding strain energy function is
\begin{align*}
  2W =& C_{11} \overline{\epsilon}_{11}^2 + 2C_{12} \overline{\epsilon}_{11} \overline{\epsilon}_{22} + 2C_{13} \overline{\epsilon}_{11} \overline{\epsilon}_{33} + 2C_{14} \overline{\epsilon}_{11} \overline{\epsilon}_{12} + 2C_{15} \overline{\epsilon}_{11} \overline{\epsilon}_{23} \\
      &+ 2C_{16} \overline{\epsilon}_{11} \overline{\epsilon}_{31} + C_{22} \overline{\epsilon}_{22}^2 + 2C_{23} \overline{\epsilon}_{22} \overline{\epsilon}_{33} + 2C_{24} \overline{\epsilon}_{22} \overline{\epsilon}_{12} + 2C_{25} \overline{\epsilon}_{22} \overline{\epsilon}_{23} \\
      &+ 2C_{26} \overline{\epsilon}_{22} \overline{\epsilon}_{31} + C_{33} \overline{\epsilon}_{33}^2 + 2C_{34} \overline{\epsilon}_{33} + \overline{\epsilon}_{12} + 2C_{35} \overline{\epsilon}_{33} \overline{\epsilon}_{23} + 2C_{36} \overline{\epsilon}_{33} \overline{\epsilon}_{31} \\
      &+ C_{44}\overline{\epsilon}_{12}^2 + 2C_{45}\overline{\epsilon}_{12} \overline{\epsilon}_{23} + 2C_{46} \overline{\epsilon}_{12} \overline{\epsilon}_{31} \\
      &+ C_{55} \overline{\epsilon}_{23}^2 + 2C_{56} \overline{\epsilon}_{23} \overline{\epsilon}_{31} \\
      &+ C_{66} \overline{\epsilon}_{31}^2
.\end{align*}

The preceding characterization is the most general with which we need to be concerned; such a material is termed \textit{anisotropic}.
Fortunately, most engineering materials possess propertied of symmetry about one or more of the planes or axes which allow the number of independent constants to be reduced.

This can be reduced for each plane of symmetry for the material ultimately leading to an isotropic material for three-axed symmetry. 

The stress-strain relationship can be written in terms of the tensor components $\epsilon_{ij}$ using the Lamé constants $\lambda$ and $\mu$:
\begin{equation*}
  \left\{ \begin{array}{cccccc}
    \sigma_{11}\\
  \sigma_{22}\\
  \sigma_{33}\\
  \sigma_{12}\\
  \sigma_{23}\\
  \sigma_{31}\\
  \end{array} \right\} = \begin{bmatrix}
    2\mu + \lambda & \lambda & \lambda & 0 & 0 & 0\\
  \lambda  & 2 \mu + \lambda & \lambda & 0 & 0 & 0\\
  \lambda & \lambda & 2 \mu + \lambda & 0 & 0 & 0\\
   &  &  & 2\mu & 0 & 0\\
   &  &  &  & 2\mu & 0\\
   &  &  &  &  &2\mu \\
  \end{bmatrix} \left\{ \begin{array}{cccccc}
    \epsilon_{11}\\
  \epsilon_{22}\\
  \epsilon_{33}\\
  \epsilon_{12}\\
  \epsilon_{23}\\
  \epsilon_{31}\\
  \end{array} \right\}
\end{equation*}
or, in indicial form,
\begin{equation*}
  \sigma_{ij} = 2\mu \epsilon_{ij} + \lambda \delta_{ij} \epsilon_{kk}
.\end{equation*}
This can also be inverted to express the strains in terms of the stresses:
\begin{equation*}
  \epsilon_{ij} = \frac{1}{2\mu}\sigma_{ij} - \frac{\lambda}{2\mu \left( 2 \mu + 3\lambda \right)} \delta_{ij} \sigma_{kk}
.\end{equation*}

Although the Lamé constants are perfectly suitable from a mathematical standpoint, it is often more convenient to use engineering material constants that are related to measurements from elementary tests.

For a uniaxial stress with $\sigma_{11}$ constant and all other $\sigma_{ij} = 0$,
\begin{align*}
  \sigma_{11} &= \left( 2\mu + \lambda \right)\epsilon_{11} + \lambda \epsilon_{22} + \lambda\epsilon_{33}, \\
  0 &= \lambda \epsilon_{11} + \left( 2\mu + \lambda \right)\epsilon_{22} + \lambda \epsilon_{33}, \\
  0 &= \lambda\epsilon_{11} + \lambda\epsilon_{22} + \left( 2\mu + \lambda \right)\epsilon_{33}
.\end{align*}
These can be equated to yield the basic form of Hooke's law:
\begin{equation*}
  \sigma_{11} = \frac{\mu \left( 2 \mu + 3\lambda \right)}{\mu + \lambda} \epsilon_{11} = E \epsilon_{11}
.\end{equation*}
From the same uniaxial stress state, the \textit{fractional contraction} may be computed as
\begin{equation*}
  - \frac{\epsilon_{22}}{\epsilon_{11}} = - \frac{\epsilon_{33}}{\epsilon_{11}} = \frac{\lambda}{2 \left( \mu + \lambda \right)} = \nu
\end{equation*}
where $\nu$ is Poisson's ratio.

A third engineering constant is obtained from the state of pure shear in two dimensions, given by $\sigma_{12} = \sigma_{21} = \text{constant}$, all other $\sigma_{ij} = 0$
\begin{equation*}
  \sigma_{12} = 2\mu \epsilon_{12} = 2 G \epsilon_{12}
\end{equation*}
where $G$ is the shear modulus.

Although $E$, $\nu$, and $G$ aer convenient, we realize only two of these are independent since $G = \mu$ and both $E$ and $\nu$ are defined in terms of $\lambda$ and $\mu$.
The relationships between the Lamé and engineering constants are collected as:
\begin{gather*}
  \mu = \frac{E}{2 \left( 1 + \nu \right)} \\
  \lambda = \frac{\nu E}{\left( 1 + \nu \right) \left( 1 - 2\nu \right)} \\
  E = \frac{\mu \left( 2\mu + 3\lambda \right)}{\mu + \lambda} \\
  \nu = \frac{\lambda}{2 \left( \mu + \lambda \right)} \\
  G = \mu
.\end{gather*}
Another relationship between the constants is defined as the bulk modulus
\begin{equation*}
  K = \frac{E}{3 \left( 1 + 2\nu \right)}
.\end{equation*}

Note that for $\lambda$ to remain finite, Poisson's ratio must lie between
\begin{equation*}
  -1 < \nu < \num{0,5} 
.\end{equation*}
Also, for $\nu = \num{0,5}$, the upper limit, the dilation $\Delta = 0$, and the material is said to be incompressible.

The explicit form of the generalized Hooke's law in terms of the engineering constants and tensoral strains is best written as
\begin{equation*}
  \Vec{\epsilon} = \Vec{C}^{-1} \Vec{\sigma}
\end{equation*}
or
\begin{equation*}
  \left\{ \begin{array}{cccccc}
    \epsilon_{11}\\
  \epsilon_{22}\\
  \epsilon_{33}\\
  \epsilon_{12}\\
  \epsilon_{23}\\
  \epsilon_{31}\\
  \end{array} \right\} = \begin{bmatrix}
    \frac{1}{E} & - \frac{\nu}{E} & - \frac{\nu}{E} & 0 & 0 & 0\\
  - \frac{\nu}{E} & \frac{1}{E} & - \frac{\nu}{E} & 0 & 0 & 0\\
  -\frac{\nu}{E} & -\frac{\nu}{E} & \frac{1}{E} & 0 & 0 & 0\\
  0 & 0 & 0 & \frac{1}{2G} & 0 & 0\\
  0 & 0 & 0 & 0 & \frac{1}{2G} & 0\\
  0 & 0 & 0 & 0 & 0 & \frac{1}{2G}\\
  \end{bmatrix} \left\{ \begin{array}{cccccc}
    \sigma_{11}\\
  \sigma_{22}\\
  \sigma_{33}\\
  \sigma_{12}\\
  \sigma_{23}\\
  \sigma_{31}\\
  \end{array} \right\}
.\end{equation*}

The elements of $\Vec{C}^{-1}$ are called \textit{compliances}.

\subsection{Transverse Isotropic Material}

In \Cref{fig:f4_5a}, a configuration with the fibers oriented along the $z$-axis, perpendicular to the symmetric $x-y$-plane is shown.
These fibers are able to stretch in a single direction as shown on \Cref{fig:f4_5b}.
This might correspond to the muscle fibers inside a muscle.

\begin{figure}
     \centering
     \begin{subfigure}[ht]{0.48\textwidth}
         \centering
         \includegraphics[width=\textwidth]{f4_5a}
         \caption{Transverse isotropic material}
         \label{fig:f4_5a}
     \end{subfigure}
     \hfill
     \begin{subfigure}[ht]{0.48\textwidth}
         \centering
         \includegraphics[width=\textwidth]{f4_5b}
         \caption{Bundled fiber configuration}
         \label{fig:f4_5b}
     \end{subfigure}
        \caption{Transverse isotropic material and bundled fiber configurations}
        \label{fig:f4_5}
\end{figure}

This is a transverse isotropic material with five elastic constants, $\lambda_{\perp}, \lambda_{\parallel}, \lambda_M, \mu_{\perp},$ and $\mu_{\parallel}$.
The $\lambda$ terms are related to the propagation of longitudinal waves in the different directions, while the $\mu$ terms are shear moduli.
The expanded stress-strain matrix $\Vec{C}_{TI}$ is:
\begin{equation*}
  \Vec{C}_{TI} = \begin{bmatrix}
    2\mu_{\parallel} + \lambda_{\parallel}& \lambda_M & \lambda_M & 0 & 0 & 0\\
  \lambda_M & 2 \mu_{\perp} + \lambda_{\perp} & \lambda_{\perp} & 0 & 0 & 0\\
  \lambda_M & \lambda_{\perp} & 2\mu_{\perp} + \lambda_{\perp} & 0 & 0 & 0\\
  0 & 0 & 0 & \mu_{\perp} & 0 & 0\\
  0 & 0 & 0 & 0 & \mu_{\parallel} & 0\\
  0 & 0 & 0 & 0 & 0 & \mu_{\parallel}\\
  \end{bmatrix}
.\end{equation*}
Note, that this is expressed in terms of Lamé constants $\mu$ and $\lambda$.
An alternate form using engineering constants is
\begin{equation*}
  \begin{bmatrix}
    \epsilon_{11}\\
  \epsilon_{22}\\
  \epsilon_{33}\\
  2\epsilon_{23}\\
  2\epsilon_{13}\\
  \epsilon_{12}\\
  \end{bmatrix} = \begin{bmatrix}
    \frac{1}{E_t} &  - \frac{\nu_{pt}}{E_{p}} & - \frac{\nu_{pt}}{E_p} & 0 & 0 & 0\\
  - \frac{\nu_{tp}}{E_t} & \frac{1}{E_p} & - \frac{\nu_p}{E_p} & 0 & 0 & 0\\
  - \frac{\nu_{tp}}{E_t} & - \frac{\nu_p}{E_p} & \frac{1}{E_p} & 0 & 0 & 0\\
  0 & 0 & 0 & \frac{1}{\mu_p} & 0 & 0\\
  0 & 0 & 0 & 0 & \frac{1}{\mu_t} & 0\\
  0 & 0 & 0 & 0 & 0 & \frac{1}{\mu_t}\\
  \end{bmatrix} \begin{bmatrix}
    \sigma_{11}\\
  \sigma_{22}\\
  \sigma_{33}\\
  \sigma_{23}\\
  \sigma_{13}\\
  \sigma_{12}\\
  \end{bmatrix}
.\end{equation*}

The seven variables in the above compliance matrix are related by the additional equations
\begin{gather*}
  \frac{\nu_{pt}}{E_p} = \frac{\nu_{tp}}{E_t} \iff \frac{E_t}{E_p} = \frac{\nu_{tp}}{\nu_{pt}} \\
  \mu_p = \frac{E_p}{2 \left( 1 + \nu_p \right)}
.\end{gather*}
Also, as the compliance matrix should be positive definite, the following five relationships exist for stability
\begin{align*}
  E_t, E_p, u_t, u_p &> 0 \\
  \left| \nu_p \right| &< 1 \\
  \nu_{pt} &< \sqrt{\frac{E_p}{E_t}} \\
  \nu_{tp} &< \sqrt{\frac{E_t}{E_p}} \\
  1 - \nu_p^2 - 2\nu_{pt}\nu_{tp} - 2\nu_{pt}\nu_{tp}\nu_p &> 0
.\end{align*}
